Landlords have been largely overlooked in research on segregation, neighborhoods and housing inequality. However, they have significant agency about the type and quality of housing in a given area. Given that households often value properties of housing units more than properties of the neighborhood and that housing is unequally distributed in space, existing theories likely misrepresent how and why residential segregation emerges and neighborhoods form. In this article, I present a theoretical model of the housing market that considers both the preferences and decisions of households as well as the agency of landlords. Based on three empirically grounded assumptions (households value the quality of a housing unit and the status of the neighborhood while landlords invest in their units when rents in the neighborhood rise), I can unify and explain multiple phenomena in urban sociology, such as the emergence and stability of segregated neighborhoods and patterns of inequalities in residential mobility, housing and neighborhood quality, as well as housing affordability.